% !TEX program = xelatex
\documentclass[14pt, a4paper, oneside]{memoir}

\input{./packages.tex}
\input{./styles.tex}
\input{./commands.tex}

\title{Обоснование стратегии оптимизации правдоподобия при наблюдаемых данных}
\begin{document}

\begin{titlingpage}
	\maketitle
\end{titlingpage}

\section{Правдоподобие смеси}

Допустим мы имеем независимые наблюдения $x_1, \dots, x_n$ полученные из смеси распределений.

Плотность смеси распределений имеет следующий вид:

\[f(x ~|~ \Psi) = \sum_{i=1}^k \omega_i \cdot f_i (x~|~\theta_i)\]

где $\Psi$ --- вектор параметров смеси $\omega_i, \theta_i$, $f_i$ --- $i$-ая компонента смеси.

Правдоподобием смеси распределений при наблюдаемых данных будет называться функция:

\[\mathcal{L}(x ~|~ \Psi) = \prod_{i=1}^n f(x_i ~|~ \Psi) = \prod_{i=1}^n \left(\sum_{j=1}^k \omega_j \cdot f_j(x_i ~|~ \theta_j)\right)\]

Для оценки параметров достаточно максимизировать правдоподобие. В силу монотонности логарифма можно оптимизировать логарифм правдоподобия:

\[\ln \mathcal{L}(x ~|~ \Psi) = \sum_{i=1}^n \ln \left(\sum_{j=1}^k \omega_j \cdot f_j(x_i~|~ \theta_j)\right)\]

\section{Покомпонентная оптимизация логарифма правдоподобия при наблюдаемых данных}

Пусть мы оптимизируем параметры $\theta_q$. Найдем градиент логарифма правдоподобия:

\[\nabla_{\theta_q} \ln \mathcal{L}(x ~|~ \Psi) = \sum_{i=1}^n \nabla_{\theta_q} \ln \left( \sum_{j=1}^k \omega_j \cdot f_j(x_i~|~ \theta_j)\right)\]

По правилу дифференцирования сложной функции:

\[\sum_{i=1}^n \nabla_{\theta_q} \ln \left( \sum_{j=1}^k \omega_j \cdot f_j(x_i~|~ \theta_j)\right) = \sum_{i=1}^n \frac{1}{\sum_{j=1}^k \omega_j \cdot f_j(x_i~|~ \theta_j)} \cdot \nabla_{\theta_q} \left(\sum_{j=1}^k \omega_j \cdot f_j(x_i~|~ \theta_j) \right)\]

В последнем сомножители все слагаемые кроме $\omega_j \cdot f_q(x_i~|~ \theta_q)$ не зависят от $\theta_q$, поэтому:

\[\nabla_{\theta_q} \left(\sum_{j=1}^k \omega_j \cdot f_j(x_i~|~ \theta_j) \right) = \omega_q \cdot \nabla_{\theta_q} f_q(x_i ~|~ \theta_q)\]

Подставим обратно:

\[\nabla_{\theta_q} \ln \mathcal{L}(x ~|~ \Psi) = \sum_{i=1}^n \frac{\omega_q \cdot \nabla_{\theta_q} f_q(x_i ~|~ \theta_q)}{\sum_{j=1}^k\omega_j \cdot f_j(x_i~|~ \theta_j)}\]

Итоговое условие экстремума:

\[\sum_{i=1}^n \frac{\omega_q \cdot \nabla_{\theta_q} f_j(x_i ~|~\theta_q)}{C + \omega_q \cdot f_q(x_i~|~ \theta_q)} = 0\]

где $C$ -- вклад остальных компонент.

\end{document}
